\documentclass{jsarticle}
\usepackage{amsmath}
\usepackage{amssymb}
\begin{document}
\title{}
\author{}
\date{}
\maketitle

       ブライアン・グリーン教授の宇宙からの演繹された説明で、
    場の理論となるオイラーのβ関数を背景にして、
    ゼータ関数が単体となり、それを元に構築された素粒子が
    Jones多項式として、各方程式の原本になり、
    それと同型の特殊相対性理論の多様体積分が、
    一般相対性理論の多様体積分として、ガンマ関数の大域的積分多様体
    として、オイラーの判別式のノルムのD-braneの種数として、
    M理論としてのカラビ・ヤウ多様体が、AdS5多様体と同型と言えるのを、
    始めは、特殊相対性理論と一般相対性理論の説明をされて、
    それから、量子力学の不確定性原理を述べて、
    そのデータが、オイラーのβ関数を場の理論とすると、
    全部の重力場と量子力学レベルの素粒子方程式が、
    同型となる、Jones多項式へと連想される説明で、
    相棒のリサ・ランドール博士と本当に友人なのですね、
    との感想で、このthe elegant Universeのの本を読んでの
    感想です。
    表紙が、種数が3の多様体を描いていて、サーストン・ペレルマン多様体
    の1部分空間を示していて、この種数3が$E^{0}\times S^{2} $と
    長生きする宇宙と、エロさまで表している、ブライアン・グリーン教授の大人の
    表現に、いいなと思いました。

    11次元多様体は、10次元で重力場を表して、11次元目でディラトンが
    出てきて、それが、AdS5多様体では、4次元目で重力場と反重力場を
    表していて、5次元目に電磁場を表していて、
    この11次元と5次元が、種数で同型と言えて、両方が、M理論を
    別の側面から見ているとリサ・ランドール博士とブライアン・グリーン教授
    は、過去も未来も現在も述べている。

\end{document}
